\section{Discussion}\label{sec:discussion}

\paragraph{A precision-oriented criterion.}
The results of Sections~\ref{sec:simulation} and~\ref{sec:application} draw a
consistent picture across synthetic fields, hyperspectral scenes and
multispectral imagery: standard internal criteria are powerful at recovering
well-separated classes but cannot recognise their absence, whereas ESS-Dip
controls the false-split rate at the cost of recall when classes are many or
weakly separated. The two behaviours are not interchangeable failures of
calibration but the two ends of a deliberate trade-off. By referring the dip
statistic to the effective rather than the nominal sample size, the test is
made conservative in exact proportion to the information the autocorrelation
removes; the level $\alpha$ and the range $\ell$ are the knobs that set where on
the precision--recall curve the method sits. We have fixed them once, without
per-scene tuning, and the resulting operating point is one of high precision:
ESS-Dip answers reliably whether discrete structure is present \emph{beyond}
spatial autocorrelation, and is correspondingly cautious about enumerating
structure that the data do not independently support. For the analyst this is
the useful direction of error, a guard against the over-segmentation that
Section~\ref{subsec:realspec} documents for the standard indices, but it is
a genuine limitation when an exhaustive class inventory is the goal. This fixes
the method's natural role: not a stand-alone classifier (on real multi-class
scenes it under-reports), but a calibrated homogeneity test, a guard placed
before an unsupervised classification or a stopping rule for a region-growing
or regionalization procedure (Section~\ref{subsec:bg-constrained}), where a
negative one can trust is worth more than an exhaustive but unreliable class
count.

\begin{revblock}
\paragraph{Choosing within the family, and embedding it in practice.}
The two calibrations serve different users. The default ESS-Dip is the choice
when a false split is the costly error (a homogeneity guard before
classification, a stopping rule, a confirmatory analysis) and whenever the
within-class marginal may be non-Gaussian, since the uniform null is
distribution-free. ESS-Dip-R is the better general-purpose class enumerator
when approximate within-class Gaussianity is plausible, as it typically is
after per-band standardisation and principal-component reduction: it recovers
substantially more structure at the same measured specificity (eleven against
five classes on Salinas; balanced score $0.868$ against $0.771$). The
imbalance sweep of Section~\ref{subsec:imbalance} makes the boundary
quantitative: minority classes below about a fifth of the scene call for the
recall variant, and below a twentieth lie beyond both. Neither variant needs a
pre-selected univariate feature: the projection is part of the procedure,
computed in the full feature space, and what the method does require is a
space in which Euclidean geometry is meaningful, so bands should be
standardised and strongly collinear cubes reduced by principal components
(we use ten throughout; Section~\ref{subsec:realspec} shows the results are
insensitive to truncations from five to thirty). Three embedding patterns
cover the practical uses: as a pre-test that decides whether any unsupervised
classification of a scene is warranted at all; as the stopping rule of a
divisive or region-growing procedure, including spatially constrained
clustering (Section~\ref{subsec:bg-constrained}); and as a post-hoc audit
that applies the homogeneity test within each cluster proposed by any
external algorithm, flagging the clusters the data do not support.
\end{revblock}

\paragraph{Relation to geostatistical null models.}
Assumption~\ref{ass:null} replaces the independence reference of the classical
criteria by a unimodal Gaussian random field, the continuous-field counterpart
of the spatial neutral models used in disease mapping \citep{Goovaerts2004} and
of Moran spectral randomization \citep{WagnerDray2015}. We realise it
implicitly, through the effective-sample-size calibration of the dip, rather
than by simulating reference fields. A more explicit route would calibrate the
statistic against an ensemble of sequential-Gaussian-simulation realisations
conditioned on the empirical variogram; this would relax the isotropy built
into our scalar range $\ell$ and accommodate anisotropic or nested
autocorrelation structures at the cost of the simulation effort. We see the
present scalar calibration and a full simulation-based null as two points on a
spectrum trading fidelity against computation, and the near-linear cost of the
former (Section~\ref{subsec:complexity}) as its principal practical advantage
on megapixel scenes. \revtext{The choice of null also encodes an analytical
judgement rather than a universal standard: the family's reference declares
autocorrelated continuous variation to be non-structure, which is the right
convention for class counting, but applications in which smooth sub-structure
is itself of interest (gradients within a single cover type that an analyst
wishes to subdivide, say) should not demand that their clusters of interest
defeat an autocorrelation null.}

\paragraph{Limitations.}
\revtext{Four} limitations bound the method's applicability. First, when many spectrally
overlapping classes coexist, as in the sixteen-class hyperspectral scenes, no
unsupervised criterion, ours included, approaches the nominal class count
(Section~\ref{subsec:realfull}); recovering such inventories requires
supervision or auxiliary information that an internal criterion cannot supply.
Second, strong class imbalance degrades recall: when one class dominates the
scene, as the crop class does in the Sentinel-2 image, the leading split
direction is governed by that majority and minority classes are not resolved.
Third, the false-split control of Proposition~\ref{prop:fpr} is an asymptotic
statement resting on the empirical-process hypothesis (C3) and a fixed
projection, not a finite-sample guarantee. \revtext{Fourth, the $2$-means
split direction is matched to balanced, spherical groups rather than to
density modes: when modes carry very different variances, or more than two
classes coexist at a node, the proposed direction can fail to expose the
multimodality or can cut through a class. The recursion re-tests every child,
and a poor direction costs recall rather than level
(Section~\ref{subsec:dip}); mode-seeking or projection-pursuit directions are
the natural refinement of this heuristic.}

\paragraph{Extensions.}
Each limitation suggests a direction. The recall-oriented variant ESS-Dip-R
realises one: trading the distribution-free uniform null for the model's
Gaussian one recovers more structure at the same specificity (eleven classes on
Salinas against five). Further recall could come from revisiting leaves at a
relaxed level, or testing minority modes against the residual of the majority,
to mitigate the imbalance that bounds the Sentinel-2 result. Replacing the scalar range by an anisotropic variogram model, and
the uniform dip null by conditional simulation, would extend the calibration to
non-stationary autocorrelation, although Section~\ref{subsec:aniso} shows the
isotropic estimate is already robust to directional autocorrelation up to an
eight-fold anisotropy ratio\revtext{, in the scenarios simulated there}. The same effective-sample-size
correction applies, unchanged, to spatio-temporal stacks, where the
decorrelation ``area'' becomes a space--time volume. Finally, because our
contribution is a stopping rule rather than a partitioning constraint, it
composes naturally with spatially constrained clustering
(Section~\ref{subsec:bg-constrained}): the calibrated test could decide when a
regionalization has produced as many regions as the data independently
support.
